% Generated draft: requires mathematical and source review.
\documentclass[11pt]{article}
\usepackage[margin=25mm]{geometry}
\usepackage{fontspec}
\setmainfont{Latin Modern Roman}
\usepackage{amsmath,amssymb,mathtools,unicode-math}
\setmathfont{Latin Modern Math}
\usepackage{enumitem,needspace,xurl,hyperref}
\setlength{\parindent}{0pt}
\setlength{\parskip}{5pt}
\setlength{\emergencystretch}{3em}
\newcommand{\N}{\mathbb{N}}
\newcommand{\Z}{\mathbb{Z}}
\newcommand{\Q}{\mathbb{Q}}
\newcommand{\R}{\mathbb{R}}
\newcommand{\C}{\mathbb{C}}
\newcommand{\F}{\mathbb{F}}
\newcommand{\A}{\mathbb{A}}
\newcommand{\PP}{\mathbb{P}}
\newcommand{\E}{\mathbb{E}}
\newcommand{\eps}{\varepsilon}
\newcommand{\dd}{\,\mathrm{d}}
\begin{document}
\section*{UnsolvedMath ID: 800534}
\textbf{Source problem code:} catalogue entry
\subsection{Definitions and mathematical statement}
Let \(\R^2\) denote the real Euclidean plane. For distinct points \(p,q\in\R^2\), define their open line segment by
\[
(p,q)=\{(1-t)p+tq:0<t<1\}.
\]
For a point set \(P\subseteq\R^2\), two distinct points \(p,q\in P\) are \emph{visible relative to \(P\)} if \((p,q)\cap P=\varnothing\). A subset \(S\subseteq P\) is \emph{mutually visible relative to \(P\)} if every two distinct points of \(S\) are visible relative to \(P\). Points are collinear if they lie on a common affine line.

With this standard interpretation of visibility, the supplied conjecture asserts that for every pair of integers \(k\geq2\) and \(\ell\geq2\), there exists a positive integer \(h(k,\ell)\) such that every point set \(P\subseteq\R^2\) with \(|P|\geq h(k,\ell)\) satisfies at least one of the following:
\begin{enumerate}
\item There exist \(\ell\) distinct points of \(P\) lying on a common affine line.
\item There exists a subset \(S\subseteq P\) with \(|S|=k\) such that
\[
\forall p,q\in S,\qquad p\ne q\ \Longrightarrow\ (p,q)\cap P=\varnothing.
\]
\end{enumerate}
The threshold depends only on \(k\) and \(\ell\), and visibility is tested against the entire set \(P\).

\textit{Unresolved scope qualification.} The supplied statement does not explicitly restrict \(P\) to be finite. The quantifier above retains its literal unrestricted wording, which includes infinite point sets. Whether the intended source formulation instead quantifies only over finite point sets requires confirmation. The supplied excerpt also does not explicitly define visibility or specify the underlying plane; the definitions above state the standard real Euclidean interpretation used in this draft.
\subsection{Short English statement}
For each \(k,\ell\geq2\), the conjecture asks for a threshold beyond which every planar point set contains either \(\ell\) collinear points or \(k\) points whose connecting open segments contain no point of the whole set. The intended restriction to finite sets, if any, remains unresolved. The registry labels the problem open; this is supplied status information, not an independently verified assessment.
\subsection{Sources}
\begin{itemize}
\item Local UnsolvedMath registry, permanent ID \(800534\), exact source problem code \texttt{catalogue entry}, title ``Visibility Graphs of Points, Segments and Polygons --- Conjecture 4.'' The supplied record attributes the problem to Ghosh and Goswami and records the proposed year as 2010.
\item Ghosh and Goswami, \emph{Unsolved Problems in Visibility Graphs of Points, Segments and Polygons} (2010), Conjecture 4. Supplied source URL: \url{https://arxiv.org/abs/1012.5187}. The supplied provenance records access on 2026-10-16 and extraction by \texttt{arxiv\_latex}; these are supplied provenance claims, not independently verified access or source checking.
\item The local record was created and updated on 2026-10-16 and supplies status ``open.'' Its upstream revision is \texttt{d371ee917ac9647fbb06cd6fc529deec6f3dda82}. This status is supplied dated registry information, not an independent literature assessment.
\end{itemize}
\end{document}
